\section{Instantiating the specification: a three-PR feasibility mapping}
\label{app:records}
An instance begins by declaring its population, criterion, horizon, and change boundary. It then assigns actors to roles, defines dependencies and visibility, and supplies the initial arrangement and improvement procedure. Evidence references identify versioned artifacts rather than undifferentiated transcript text. Each patch records the state it expects and the evidence its proposer could inspect. A compact contract can use the following interchange form; the referenced artifacts contain the population, protocol, scores, and retention details.
\begin{lstlisting}
{
  "id": "coverage-change-1",
  "target": "evaluation.coverage_program",
  "expected_version": 0,
  "actor": "review_board",
  "transformation": {"from": "Biased", "to": "Balanced"},
  "evidence": ["program-trials-1.v1"],
  "comparison": "population-and-net-value-criterion.v1",
  "cost_ledger": "labels-outputs-transition-1.v1",
  "horizon_rounds": 8,
  "retention": "next-review-rule.v1"
}
\end{lstlisting}
This example is a proposed full contract. The executable \texttt{Patch} class implements target, version, actor, evidence, reversibility, and admission cost; other contract fields are supplied by the surrounding experiment protocol. Reversibility is recorded but the checker does not simulate rollback. A diagnostic coverage trace contains per-round sample counts, estimates, decisions, cost totals, and authorized program changes. The outcome equations reconstruct net value from the selected workflow and ledger, while negative fixtures exercise admission failures.

The public mapping uses a convenience-sample snapshot from \texttt{pallets/flask}: the first three merged PRs before September 30, 2026 among the first 100 closed PRs returned in descending creation order. The snapshot was retrieved on September 28, 2026 at 17:45:43 UTC. Selection candidates, URLs, timestamps, and response hashes are retained. The snapshot contains factual identifiers and states, excluding message bodies and email addresses. Released actor identifiers are stable pseudonyms; account IDs and usernames are omitted. PR links remain available for provenance, so this is pseudonymization rather than irreversible anonymization.

For each PR the adapter retrieves commits, reviews, and checks at its recorded head SHA, checking pagination for the subordinate lists. \Cref{tab:public} summarizes the mapped records. Commits are artifact versions; opening, review, check, and merge records are events. Endpoint definitions support these mappings \citep{githubreviews,githubchecks}. Earlier heads, test-merge commits, off-platform activity, and historical branch protection require additional sources.
\begin{table}[!htbp]
\caption{Public-record snapshot. These counts describe record availability, not organizational improvement or AI participation.}
\label{tab:public}\centering\small
\begin{tabular}{@{}lrrrr@{}}\toprule
PR & Commit artifacts & Reviews & Head checks & Events \\\midrule
\href{https://github.com/pallets/flask/pull/6133}{\#6133} & 1 & 0 & 14 & 16 \\
\href{https://github.com/pallets/flask/pull/6096}{\#6096} & 2 & 5 & 14 & 21 \\
\href{https://github.com/pallets/flask/pull/6095}{\#6095} & 1 & 3 & 14 & 19 \\\bottomrule
\end{tabular}
\end{table}
Historical permission and actual evidence exposure remain unknown even when a merge and a public review are visible. The adapter preserves these unknown fields and does not pass the records to the complete-trace validator. No workflow intervention or counterfactual is observed. A generic log augmented with the same missing information would support the same checks; the practical benefit tested here is identifying the required connections and missing instrumentation.

\section{Analytical background for interaction and evaluation}
\label{app:analytical}
\subsection{Coverage and indistinguishability}
Suppose two environments induce identical distributions over every history accessible to a procedure, but a specified workflow change improves the population in one and harms it in the other. Any binary sign decision based on those histories has the same output distribution in both environments. If its probability of reporting improvement is $q$, its two error probabilities are $1-q$ and $q$. Their sum is one, so at least one is at least $1/2$. This is the standard observation-equivalence argument.

A construction uses equally frequent flagged and unflagged cases. A change lowers flagged error from $.20$ to $.10$ in both environments; unflagged error falls to $.10$ in one and rises to $.50$ in the other. If only flagged labels are observed and no downstream signal arrives, evidence is identical although population risk changes from $.20$ to $.10$ or $.30$. With known positive inclusion probabilities $q_i$ and accurate labels $E_i$, the Horvitz--Thompson estimator $N^{-1}\sum_i Z_i E_i/q_i$ is design-unbiased for finite-population risk \citep{ht}. Small inclusion probabilities can still produce high variance. Study 1 has positive coverage in every cell and tests this finite-evidence problem; it is not an instance of exact observational equivalence.

\subsection{Dependent evidence and interaction order}
For equal-variance errors with equicorrelation $\rho$, the variance of the mean is $\sigma^2[\rho+(1-\rho)/n]$, with $-1/(n-1)\leq\rho\leq1$. Common bias contributes an additional squared-bias term to mean squared error. Information dependence therefore limits the gain from adding actors \citep{clemen}; more actors alone need not produce proportionally more information.

Let independent initial human and agent estimation errors have variances $v_H$ and $v_A$. Suppose the human sees the agent answer and replaces the private judgment by $(1-\alpha)H+\alpha A$. Equal averaging of that judgment with $A$ has risk
\begin{equation}
R(\alpha)=\{(1-\alpha)^2v_H+(1+\alpha)^2v_A\}/4.
\end{equation}
For $v_A<v_H$ and $0\leq\alpha\leq1$, the optimum is $\alpha^*=(v_H-v_A)/(v_H+v_A)$, giving $v_Hv_A/(v_H+v_A)$. This is also the risk of optimal inverse-variance weighting of the original private judgments. Thus a gain over equal private averaging is a reweighting gain, not evidence that exposure creates information. With initial covariance $\sigma_{HA}$, risk gains the term $(1-\alpha^2)\sigma_{HA}/2$. The parameters describe a candidate interaction model; they are not empirical estimates of anchoring.

Cross-functional breadth can be represented by reduced transfer cost and changed capability profiles. A simple accounting decomposition is saved handoff cost plus improved integration, minus lost specialist depth and switching cost. Its sign depends on measured magnitudes and on the dependency graph. The specification preserves these separate terms instead of assigning a universal benefit to generalists.

\section{Study 1: complete specification}
\label{app:protocol}
\subsection{Core protocol and exploratory controls}
The core protocol, \path{examples/learning_protocol.json}, records a design fixed before execution and is archived at commit \texttt{75d312e}. This first public commit contains both protocol and results; it does not establish a pre-run Git timestamp. The protocol specifies 48 rounds, a 2736-label ceiling for each eight-round block, three catalog programs, six arms, three memory conditions, three environments, and the full allocation grid. There was no parameter selection on test results. The trial-matched control and change-timing checks were specified after inspection of the core results to examine a possible acquisition-timing explanation. Their design and independent random streams are described below. They use fresh streams and are reported as exploratory attribution checks. Neither protocol is an external preregistration.

The core has $3\times3\times6\times128=6912$ trajectories. The allocation grid uses 96 replicates for each share/trial-count combination and each discovery arm under stationary cumulative memory and reversal with Window 8. Four fixed references and two memory ablations in each environment bring this analysis to 4608 trajectories. The exploratory controls cross four environments (stationary and reversals before rounds 21, 25, 29), three memory conditions, four arms, and 128 replicates: 6144 trajectories. The three independent keyed root seeds are 920000, 930000, and 940000. Each random stream is indexed by root seed, round, acquisition stage, program, trial, and replicate. Conditions share underlying potential-label streams within each root; no policy receives unqueried labels.

\subsection{Evidence update and program trials}
For a template--stratum cell, cumulative memory adds each actually queried label once. Window 8 uses observations at historical rounds $t-8$ through $t-1$, plus current observations. Reset retains no earlier observations. All program trials at a review receive the same pre-review memory snapshot. Trial outputs and fresh validation labels are pooled only after program selection, then supplied to the current operational screen. Validation observes only the selected template, while screening observes all templates except for the second phase of the rejection comparator. A label from a rejected program remains usable if its source and population are appropriate. The pooling ablation withholds all program-trial labels from workflow memory while retaining program scores.

Each trial score is $1-3\widehat r^{\rm val}-.22N^{\rm screen}/4096$, where validation uses equal counts across the two strata. A program's score is the mean of its trial means retained under the memory rule. Each review has the same trial count within a configuration. Ties favor Biased, then Balanced, then Neyman. Program-score memory can be disabled independently of workflow memory. Retained scores were generated under different historical knowledge states; they are empirical assessments used by this specified controller, not identically distributed observations of an invariant program value.

For block budget $B=2736$, program share $f$, three programs, and $q$ trials per program, a trial receives $b=\lfloor\lfloor fB\rfloor/(3q)\rfloor$ label units. Validation obtains $v=\max\{1,\lfloor\lfloor .2b\rfloor/2\rfloor\}$ labels per stratum; screening receives $b-2v$ total units. The base $(f,q)=(.2,1)$ therefore allows 182 units per trial, with 18 validation labels per stratum and 146 units for screening. Actual counts can be smaller after integer allocation. Let $M_i$ be the actual number of trial labels for replicate $i$ at a block's start. Its subsequent per-round operational allowance is $\lfloor(B-M_i)/8\rfloor$. Discover once has $M_i=0$ after block 1. This rule keeps the budget fixed when trial count or program share changes.

\subsection{Within-template and between-template allocation}
For the catalog programs, let $u$ be one template's share of the operational or trial allowance. Balanced obtains $\lfloor u/2\rfloor$ labels in each stratum. Biased obtains $\max\{1,\lfloor .95u\rfloor\}$ and $\max\{1,\lfloor .05u\rfloor\}$. All tested allocations fit the allowance. Neyman first obtains $m=\min\{2,\max\{1,\lfloor u/4\rfloor\}\}$ labels per cell. Using retained and pilot data, it estimates $p_s$, sets $v_s=p_s(1-p_s)$, and allocates
\begin{equation}
n_s=m+\left\lfloor\frac{(u-2m)\sqrt{v_s}}{\sum_\ell\sqrt{v_\ell}}\right\rfloor.
\end{equation}
Pilot labels enter the posterior once. Equal label costs make this the equal-cost stratified allocation. Catalog selection breaks posterior ties in template order standard, specialized, broad.

The successive-rejection comparator treats one pair of stratum labels as a bounded loss with mean equal to population risk. With $n$ affordable pairs and $K=3$, it uses the schedule of \citet{bai}:
\begin{equation}
\overline{\log}K=\frac12+\sum_{j=2}^{K}\frac1j,\qquad
n_k=\left\lceil\frac{n-K}{(\overline{\log}K)(K+1-k)}\right\rceil,\quad k=1,2.
\end{equation}
All three templates receive $n_1$ pairs; the worst posterior-mean template is removed, and the remaining two reach $n_2$ pairs each. Elimination ties remove the lowest indexed template; final ties use the lowest surviving index. Historical evidence enters the posterior under the memory conditions. Thus the implementation uses the published rejection schedule with declared posterior ranking and deterministic ties. At the fixed allowance of 342 labels it uses 336 labels. No theoretical bound is claimed for this modification under memory or drift. SR is a fixed between-template comparator and is not part of the three-program discovery catalog.

\subsection{Outcome accounting, checks, and uncertainty}
For deployed template $k_t$ with true population risk $r_{k_t,t}$, total queried labels $Q_t$, and program-change indicator $S_t$, the external score is
\begin{equation}
V_t=1-3r_{k_t,t}-.02-\{.22Q_t+.50S_t\}/4096.
\label{eq:net}
\end{equation}
Each queried label corresponds to a generated evaluation output, so $.22$ includes $.20$ labeling and $.02$ generation. The separate $.02$ term prices each production output. Unqueried potential-label arrays are a simulator device, not actor-visible generated outputs. All rejected-program trials are charged. Program comparison costs are not subtracted a second time after computing net value.

The simulator checks every block ceiling and the identity
$V_t=1-3r_t^*-.02-\mathrm{regret}_t-\mathrm{expense}_t$,
where $r_t^*$ is the best available risk. An independent scalar reconstruction checks 36 fixed-Balanced trajectories across the three environments and memory conditions to absolute tolerance $10^{-12}$. Separate fixtures check the lookback boundary, rejection allocation, and selection of a known best template. The patch checker enforces target authority and state version for program changes.

Reported half-widths are $1.96s/\sqrt N$ across independent replicate trajectories. Paired intervals use within-replicate differences under the common streams. They quantify simulation uncertainty conditional on the supplied parameters. The uniform fixed benchmark averages the conditional mean values of Biased, Balanced, and Neyman under an initial uniform program prior; its paired contrast integrates over this discrete prior rather than adding random program draws. No claim of equivalence or correction for multiple hypothesis testing is attached to intervals spanning zero.

\Cref{tab:base} reports all core net-value means. Supplementary Table S1 reports their Monte Carlo intervals, costs, harmful adoption, selection regret, and late outcomes; replicate-level data accompany the reproducibility materials.
\begin{table}[!htbp]
\caption{Complete core net-value means. B: Biased; Bal: Balanced; N: Neyman; Mix: uniform fixed-program mixture; SR: successive rejection; O: discover once; R: repeated discovery. Each cell uses 128 trajectories. Monte Carlo intervals and additional outcomes are reported in Supplementary Table S1.}
\label{tab:base}\centering\small
\begin{tabular}{@{}llrrrrrrr@{}}\toprule
Environment & Memory & B & Bal & N & SR & O & R & Mix \\\midrule
\tableinput{sections/learning_base_rows.tex}
\bottomrule\end{tabular}\end{table}

\Cref{tab:grid} reports the complete factorial grid, so favorable settings are not selected for the main result. Supplementary Table S2 reports retained-program proportions, conditional sample sizes, and net value and harmful adoption over the subsequent eight rounds. These conditional means remain descriptive because the retained program is selected using noisy evidence.
\begin{table}[!htbp]
\caption{Repeated discovery minus discover-once at fixed total budget. All program shares and trial counts are shown. Entries are paired differences $\pm$ 95\% Monte Carlo half-width; 96 fresh replicates per condition.}
\label{tab:grid}\centering\small
\begin{tabular}{@{}rrrr@{}}\toprule
Program share & Trials/program & Stationary, cumulative & Reversal, Window 8 \\\midrule
\tableinput{sections/learning_grid_rows.tex}
\bottomrule\end{tabular}\end{table}
\begin{table}[!htbp]
\caption{Repeated discovery: retained-program percentages and subsequent eight-round outcomes conditional on retaining Balanced. Stationary and uniform gain use cumulative memory; reversal uses Window 8. B/Bal/N are Biased/Balanced/Neyman; $n_{\rm Bal}$ is the number retaining Balanced among 128 replicates. Harm is the percentage of rounds deploying a template worse than standard. Supplementary Table S2 reports all three retained programs and all three memory conditions for both discover-once and repeated-discovery rules across the three environments.}
\label{tab:programs}\centering\small
\begin{tabular}{@{}lrrrrrrr@{}}\toprule
Environment & Review & B (\%) & Bal (\%) & N (\%) & $n_{\rm Bal}$ & Net given Bal & Harm (\%) \\\midrule
\tableinput{sections/learning_program_rows.tex}
\bottomrule\end{tabular}\end{table}
\FloatBarrier


\section{Supplementary audit-control model}
\label{app:audit}
The supplementary audit model isolates longitudinal evidence acquisition under a fixed controller. It contains 160 periods of 32 tasks. Route A has base value 1 and error probability $.04$ (stationary) or alternating $.04/.25$ blocks of 40 periods; route B has base value $.82$ and error probability $.08$. An error costs 3, an audit costs $.20$, and an audit-rate change costs $.50$ per batch. Only audits reveal A's errors, available next period; executing B can still be paired with acquiring an A audit. Observations update a sliding-window Beta estimate. The base prior is $\operatorname{Beta}(1,9)$ and window 12.

The payoff-optimal routing threshold is $.14$, since $1-3p=.82-3(.08)=.58$. The heuristic adaptive controller audits at rate $.40$ when posterior $P(p>.14)$ lies between $.10$ and $.90$, and at $.025$ otherwise. Periodic dense auditing uses $.40$ for the first three periods of each 12-period cycle and $.025$ otherwise. The base audit comparison uses 400 seeds from 20260928, common across policies. Stationary net values are $.86089$ for fixed $.025$, $.86395$ for fixed $.05$, $.85841$ for fixed $.10$, and $.87976$ for no auditing. No auditing retains A under the base prior, the optimal route in that stationary case; audits cannot repair current outputs.

With slow shifts, adaptive auditing exceeds fixed $.10$ by $.00857$ per task (paired interval $.00676$--$.01038$). A separate 200-seed grid, beginning at 20360928, crosses windows 6, 12, 24 with stationary, 40-period, and 10-period regimes. It includes no audit, fixed $.05/.10/.40$, periodic dense, and adaptive policies. Supplementary Table S3 reports all 54 combinations of environment, memory window, and audit policy.

\subsection{Finite-horizon EVSI controller}
The finite-horizon expected-value-of-sample-information (EVSI) controller evaluates acquiring $m\in\{0,1,4,16\}$ labels. Given current Beta parameters $(a,b)$, a beta-binomial predictive distribution over $k$ observed errors gives posterior mean $(a+k)/(a+b+m)$. Let $L_\theta(p)$ be $1-3p$ if $p\leq\theta$ and $.58$ otherwise. The controller maximizes
\begin{equation}
h\{\E_k[L_\theta((a+k)/(a+b+m))]-L_\theta(a/(a+b))\}
-\frac{\bigl[(.20+g_c I_B)m+.50I_{\rm change}\bigr]}{32}.
\end{equation}
$I_{\rm change}$ equals 1 when the audit rate changes and 0 otherwise. Both costs in the bracket are subtracted. $I_B$ indicates executing B; $g_c$ is the cost of generating a counterfactual A output to audit. The predictive sum is exact under the current Beta model. The controller approximates future value with horizon $h$ and a sliding window, so it is a finite-horizon EVSI heuristic under drift.

Development fixes the prior at $\operatorname{Beta}(1,9)$, the threshold at $.14$, and $g_c=0$. Seeds 710000--710063 choose $h$ from $1,4,8,16$ using stationary and slow-shift mean value; $h=16$ is selected. Testing uses 96 independent seeds 810000--810095, three priors, thresholds $.14/.16$, generation costs $g_c=0/.20$, two environments, and three policies: 6912 trajectories. Threshold $.16$ is a misspecified routing-rule sensitivity under unchanged payoffs. Nonzero $g_c$ charges selectively generated A outputs while B is executed; it is distinct from a common cost charged for producing every candidate under every policy.

\Cref{tab:audit} reports all prior/threshold means at $g_c=0$. Supplementary Table S4 reports both candidate-generation cost conditions, marginal Monte Carlo intervals, and paired policy contrasts. EVSI's performance depends strongly on the prior: with no initial feedback, information acquisition itself can stall. Four regression cases match a reference scalar implementation to $10^{-12}$.
\begin{table}[!htbp]
\caption{Supplementary audit sensitivity at zero additional candidate-generation cost. Columns give net-value means over 96 trajectories; F: fixed $.10$, A: adaptive heuristic, E: EVSI. Uncertainty estimates and results at $g_c=.20$ are reported in Supplementary Table S4.}
\label{tab:audit}\centering\small
\begin{tabular}{@{}llrrrrrr@{}}\toprule
 & & \multicolumn{3}{c}{Stationary} & \multicolumn{3}{c}{Slow shifts} \\
Prior $(a:b)$ & Threshold & F & A & E & F & A & E \\\midrule
\tableinput{sections/audit_rows.tex}
\bottomrule\end{tabular}
\end{table}

